\documentclass[a4paper,12pt]{report}
\usepackage[utf8]{inputenc}
\usepackage[T2A]{fontenc}
\usepackage[russian]{babel}
\usepackage{amsmath}
\usepackage{listings}
\usepackage{xcolor}
\usepackage{tabularx}
\usepackage{longtable}
\usepackage{hyperref}

\begin{document}

\chapter{ТЕОРЕТИЧЕСКИЕ ОСНОВЫ ПРОЕКТИРОВАНИЯ SQL-ДВИЖКА}

Прежде чем приступать к описанию архитектуры и программной реализации, необходимо заложить теоретическую базу, на которой строится всё остальное. В этой главе рассматриваются концепции, без понимания которых невозможно объяснить ни одно из принятых инженерных решений: почему выбран именно такой алгоритм разбора, почему реализованы две модели исполнения, и почему Python с NumPy является обоснованным — а не случайным — выбором инструментальной среды.

Изложение построено вокруг единого демонстрационного запроса, который проходит сквозь всю главу от начала до конца:

\begin{verbatim}
SELECT name, SUM(salary)
FROM employees
WHERE dept_id = 10
GROUP BY name
ORDER BY SUM(salary) DESC;
\end{verbatim}

К концу главы будет понятно, почему для исполнения этих пяти строк текста требуется целый программный комплекс, а не просто цикл по массиву данных.

\vspace{0.5cm}\hrule\vspace{0.5cm}

\section{SQL как декларативный язык программирования}

До начала 1970-х годов взаимодействие с базами данных носило сугубо навигационный характер. Разработчик был обязан вручную описывать алгоритм обхода данных: указывать, с какого узла начать, по какому указателю перейти дальше, когда остановиться. По сложности это напоминало программирование на языке ассемблера — формально задача решаема, но ценой колоссальных усилий и хрупкого кода.

В 1970 году сотрудник IBM Эдгар Кодд опубликовал работу «A Relational Model of Data for Large Shared Data Banks», в которой предложил принципиально иной взгляд на управление данными. Идея состояла в том, чтобы представить данные не как связанные узлы в графе, а как математические отношения — множества кортежей с именованными атрибутами. Взаимодействие с такими структурами описывается не алгоритмами обхода, а логическими предикатами: «дай мне все кортежи, для которых выполняется условие C». Исполнение этого предиката — задача системы, а не программиста.

Практической реализацией идей Кодда стал язык SEQUEL, разработанный в лаборатории IBM San Jose и впоследствии переименованный в SQL. Важно подчеркнуть: SQL с самого начала проектировался как \textbf{специализированный язык высокого уровня (Domain-Specific Language)}, а не как системная утилита. Его синтаксис, напоминающий английскую прозу, был осознанным выбором — попыткой сделать логику данных читаемой без глубокой технической подготовки.

С развитием стандартов возможности языка существенно расширились. SQL-86 и SQL-92 зафиксировали ядро реляционной алгебры. Стандарт SQL:1999 стал переломным: с введением рекурсивных общих табличных выражений (Recursive CTE) язык формально стал \textbf{полным по Тьюрингу} — то есть на нём теоретически можно выразить любой вычислимый алгоритм. Современные стандарты добавили оконные функции, поддержку JSON и временных рядов, окончательно закрепив SQL как основной язык аналитической обработки данных.

Ключевое архитектурное следствие декларативности — полная изоляция логики от реализации. В Python или C++ программист описывает \textbf{путь} вычисления (последовательность операций). В SQL программист описывает \textbf{результат} (условие, которому должны удовлетворять данные), а движок самостоятельно выбирает путь — через механизм оптимизации. Именно этот механизм является предметом данной работы.

\begin{longtable}{|l|l|l|}
\hline
Характеристика & Императивные языки (Python, C++) & Декларативный SQL \\ \hline
Единица описания & Алгоритм (шаги выполнения) & Условие (логическое выражение) \\ \hline
Выбор пути исполнения & Программист & Оптимизатор движка \\ \hline
Оптимизация & Ручная & Автоматическая \\ \hline
Тьюринг-полнота & Прямая (циклы, рекурсия) & Через Recursive CTE \\ \hline
\end{longtable}

В рамках данной работы SQL рассматривается не просто как средство выборки данных, а как \textbf{целевая вычислительная среда}: наша задача — реализовать ядро, которое транслирует декларативные конструкции в эффективный физический план выполнения.

\vspace{0.5cm}\hrule\vspace{0.5cm}

\section{Теория формальных языков и грамматик}

Чтобы построить движок, способный «понимать» SQL, необходимо сначала формализовать само понятие «правильного» SQL-запроса. Этим занимается теория формальных языков.

В 1956 году Ноам Хомский предложил иерархию, классифицирующую формальные языки по сложности порождающих грамматик и необходимых вычислительных ресурсов для их распознавания. Согласно этой классификации, существует четыре уровня:

\begin{longtable}{|l|l|l|l|}
\hline
Тип & Название & Распознаватель & Пример \\ \hline
Тип 0 & Рекурсивно-перечислимые & Машина Тьюринга & Естественные языки \\ \hline
Тип 1 & Контекстно-зависимые & Линейно-ограниченный автомат & Некоторые синтаксические структуры \\ \hline
\textbf{Тип 2} & \textbf{Контекстно-свободные} & \textbf{Магазинный автомат (стек)} & \textbf{SQL, Java, C++} \\ \hline
Тип 3 & Регулярные & Конечный автомат & E-mail, номера телефонов \\ \hline
\end{longtable}

SQL принадлежит к \textbf{Типу 2 — контекстно-свободным языкам}. Это не произвольное утверждение: в нашем демонстрационном запросе выражение \texttt{SUM(salary)} допускает произвольно глубокое вложение функций. Регулярные выражения (Тип 3) принципиально не способны отследить соответствие открывающих и закрывающих скобок при вложенности больше фиксированной глубины — для этого нужна память, то есть стек.

Формально любая контекстно-свободная грамматика $G$ задаётся как четвёрка:

\begin{equation*}
G = (V,\; T,\; P,\; S)
\end{equation*}

где $V$ — множество нетерминальных символов (таких как \texttt{SelectStmt}, \texttt{Expression}), $T$ — множество терминальных символов (конкретные ключевые слова и идентификаторы), $P$ — правила вывода (например, \texttt{expression → expression "+" expression}), $S$ — начальный символ. В реализуемом движке эта грамматика описана декларативно в файле \texttt{parser/grammar.lark}, где каждое правило $P$ соответствует одной строке файла.

Поскольку SQL — это контекстно-свободный язык, для его разбора требуется автомат со стеком (магазинный автомат). Когда парсер встречает открывающую конструкцию — скобку подзапроса или начало агрегатной функции — он помещает маркер в стек. При встрече закрывающей конструкции маркер снимается. Это обеспечивает корректную обработку любой глубины вложенности.

\vspace{0.5cm}\hrule\vspace{0.5cm}

\section{Лексический анализ}

Лексический анализ — первый этап трансляции запроса. Его задача состоит в преобразовании непрерывной строки символов в структурированный поток дискретных единиц, называемых \textbf{токенами}. Лексер, по сути, отвечает на вопрос: «какие слова содержит этот запрос, и к каким категориям они относятся?»

Необходимо разграничить три понятия, между которыми легко запутаться. \textbf{Лексема} — это конкретная последовательность символов в тексте (например, \texttt{dept\_id} или \texttt{10}). \textbf{Токен} — это структура данных, описывающая лексему: она содержит тип и атрибутивное значение. \textbf{Паттерн} — это правило, определяющее, какие лексемы порождают токен данного типа; паттерны записываются в виде регулярных выражений.

Для нашего демонстрационного запроса лексер последовательно выделяет следующие токены:

\begin{longtable}{|l|l|l|l|}
\hline
№ & Лексема & Тип токена & Примечание \\ \hline
1 & \texttt{SELECT} & \texttt{KW\_SELECT} & Ключевое слово \\ \hline
2 & \texttt{name} & \texttt{IDENTIFIER} & Имя колонки \\ \hline
3 & \texttt{,} & \texttt{COMMA} & Разделитель \\ \hline
4 & \texttt{SUM} & \texttt{IDENTIFIER} & Агрегатная функция \\ \hline
5 & \texttt{(} & \texttt{LPAR} & Открывающая скобка \\ \hline
6 & \texttt{salary} & \texttt{IDENTIFIER} & Имя колонки \\ \hline
7 & \texttt{)} & \texttt{RPAR} & Закрывающая скобка \\ \hline
8 & \texttt{FROM} & \texttt{KW\_FROM} & Ключевое слово \\ \hline
9 & \texttt{employees} & \texttt{IDENTIFIER} & Имя таблицы \\ \hline
10 & \texttt{WHERE} & \texttt{KW\_WHERE} & Ключевое слово \\ \hline
11 & \texttt{dept\_id} & \texttt{IDENTIFIER} & Имя колонки \\ \hline
12 & \texttt{=} & \texttt{EQ} & Оператор сравнения \\ \hline
13 & \texttt{10} & \texttt{NUMBER} & Числовой литерал \\ \hline
14 & \texttt{GROUP} & \texttt{KW\_GROUP} & Ключевое слово \\ \hline
15 & \texttt{BY} & \texttt{KW\_BY} & Ключевое слово \\ \hline
16 & \texttt{name} & \texttt{IDENTIFIER} & Ссылка на колонку \\ \hline
17 & \texttt{ORDER} & \texttt{KW\_ORDER} & Ключевое слово \\ \hline
18 & \texttt{BY} & \texttt{KW\_BY} & Ключевое слово \\ \hline
19–22 & \texttt{SUM(salary)} & \textit{(повтор 4–7)} & Выражение в ORDER BY \\ \hline
23 & \texttt{DESC} & \texttt{KW\_DESC} & Направление сортировки \\ \hline
24 & \texttt{;} & \texttt{SEMICOLON} & Конец оператора \\ \hline
\end{longtable}

91 символ исходного текста трансформировался в 24 токена. Именно этот поток, а не строка, передаётся на следующий этап.

В основе работы лексера лежит механизм \textbf{регулярных выражений}, каждое из которых математически соответствует \textbf{детерминированному конечному автомату (ДКА)}. Лексер итерирует входную строку посимвольно: текущий символ определяет переход между состояниями автомата. Достигнув принимающего состояния в точке, откуда продолжить движение невозможно, автомат фиксирует лексему. При конкуренции двух правил применяется принцип \textbf{Longest Match} — побеждает правило, поглощающее больше символов. Именно поэтому \texttt{GROUP} читается как один токен, а не \texttt{GRO} + \texttt{UP}.

В реализуемом движке леческая спецификация описана декларативно в файле \texttt{parser/grammar.lark}:

\begin{verbatim}
SELECT:     "SELECT"i
FROM:       "FROM"i
WHERE:      "WHERE"i
GROUP:      "GROUP"i
BY:         "BY"i
IDENTIFIER: /[a-zA-Z_][a-zA-Z0-9_]*/
NUMBER:     /-?\d+(\.\d+)?/
%ignore     /\s+/
\end{verbatim}

Директива \texttt{%ignore} автоматически отфильтровывает пробельные символы — «шум», необходимый человеку для читаемости, но бессмысленный для машины.

Помимо токенизации, лексер выполняет ещё одну оптимизацию: повторяющиеся идентификаторы (в нашем запросе \texttt{name} встречается дважды — в позициях 2 и 16) сохраняются в \textbf{таблице символов} как числовые индексы. Вместо двух строковых объектов в памяти хранится один, и везде используется ссылка на него. Сравнение целых чисел выполняется значительно быстрее, чем посимвольное сравнение строк — что ощутимо при обработке сложных запросов с многократно повторяющимися именами.

\vspace{0.5cm}\hrule\vspace{0.5cm}

\section{Синтаксический анализ}

Синтаксический анализ, или парсинг, решает задачу более высокого уровня: он выясняет, образует ли последовательность токенов структуру, предусмотренную грамматикой языка, и строит иерархическое дерево, отражающее вложенность конструкций запроса.

В данном проекте используется алгоритм \textbf{LALR(1)} — Look-Ahead Left-to-Right с одним токеном предпросмотра. В отличие от нисходящих LL-парсеров, LALR работает \textbf{снизу вверх}: он накапливает токены в стеке и при распознавании полной правой части грамматического правила выполняет свёртку — заменяет накопленные токены нетерминальным символом из левой части. Два ключевых действия парсера — \textbf{перенос} (shift, добавить токен в стек) и \textbf{свёртка} (reduce, заменить вершину стека нетерминалом) — управляются таблицей переходов, которую Lark строит автоматически по файлу грамматики.

Префикс «LA» означает, что перед принятием решения парсер заглядывает на один токен вперёд. Это делает процесс полностью детерминированным: в каждом состоянии существует ровно одно возможное действие. Если в грамматике возникает неоднозначность (конфликт «перенос-свёртка»), это сигнализирует об ошибке проектирования грамматики, которую необходимо устранить.

Результатом работы парсера является \textbf{Abstract Syntax Tree (AST)} — дерево, в узлах которого находятся смысловые конструкции запроса, а не грамматические правила. Для нашего демонстрационного запроса AST выглядит следующим образом:

\begin{verbatim}
SelectNode
├── columns:
│   ├── ColumnRef("name")
│   └── AggFunc("SUM", arg=ColumnRef("salary"))
├── from_table: TableRef("employees")
├── where:
│   └── BinaryOp("=", ColumnRef("dept_id"), Literal(10))
├── group_by:
│   └── ColumnRef("name")
└── order_by:
    └── OrderItem(AggFunc("SUM", ColumnRef("salary")), direction="DESC")
\end{verbatim}

Это дерево является полным и однозначным представлением семантики запроса. Начиная с данного момента, строковый текст SQL движку более не нужен: все последующие этапы — планирование, оптимизация, исполнение — работают исключительно с этой структурой данных.

Следует уточнить, что Lark первично строит \textbf{Concrete Syntax Tree (CST)} — дерево разбора, сохраняющее все промежуточные правила грамматики и служебные символы. Оно избыточно для исполнения. Модуль \texttt{parser/ast\_builder.py} трансформирует CST в компактный AST: многоуровневая цепочка \texttt{select\_stmt → select\_clause → select\_item → identifier → "name"} сворачивается в один узел \texttt{ColumnRef("name")}. Такое разделение делает AST независимым от конкретной грамматики: изменения в \texttt{grammar.lark} не затрагивают описания узлов в \texttt{ast\_nodes/}.

\vspace{0.5cm}\hrule\vspace{0.5cm}

\section{Реляционная алгебра}

Построенного AST достаточно, чтобы понять \textbf{что} требуется вычислить. Однако для того, чтобы рассуждать о \textbf{корректности} преобразований запроса — например, доказать, что переставленный фильтр даёт тот же результат — необходим формальный математический язык. Таким языком является реляционная алгебра.

Реляционная алгебра располагает набором операций над отношениями (таблицами):

- \textbf{Селекция} $\sigma_{c}(R)$ — выборка кортежей из $R$, удовлетворяющих предикату $c$.
- \textbf{Проекция} $\pi_{L}(R)$ — оставить только атрибуты из списка $L$.
- \textbf{Декартово произведение} $R \times S$ — все комбинации кортежей двух отношений.
- \textbf{Тета-соединение} $R \bowtie_{\theta} S$ — эквивалентно $\sigma_{\theta}(R \times S)$.
- \textbf{Группировка и агрегация} $\gamma_{G;\, F}(R)$ — разбивка по группам $G$ с применением агрегатов $F$.

Трансформируем наш демонстрационный запрос в выражение реляционной алгебры. Именно это выражение строит логический планировщик (\texttt{planner/logical\_operators.py}) на основе AST:

\begin{equation*}
\pi_{\text{name},\, \text{SUM(salary)}} \Bigl( \tau_{\text{SUM(salary) DESC}} \bigl( \gamma_{\text{name};\, \text{SUM(salary)}} \bigl( \sigma_{\text{dept\_id} = 10}(\text{employees}) \bigr) \bigr) \Bigr)
\end{equation*}

Читая выражение изнутри наружу, мы получаем точный порядок операций:

\begin{longtable}{|l|l|l|}
\hline
Уровень & Операция & SQL-конструкция \\ \hline
Внутренний & $\sigma_{\text{dept\_id}=10}(\text{employees})$ & \texttt{WHERE dept\_id = 10} \\ \hline
2-й & $\gamma_{\text{name};\, \text{SUM(salary)}}(\ldots)$ & \texttt{GROUP BY name} + \texttt{SUM(salary)} \\ \hline
3-й & $\tau_{\text{SUM(salary) DESC}}(\ldots)$ & \texttt{ORDER BY SUM(salary) DESC} \\ \hline
Внешний & $\pi_{\text{name},\, \text{SUM(salary)}}(\ldots)$ & \texttt{SELECT name, SUM(salary)} \\ \hline
\end{longtable}

Самое важное свойство реляционной алгебры — существование \textbf{законов эквивалентности}, гарантирующих, что различные выражения дают одинаковый результат. Два закона являются основой нашего оптимизатора:

\textbf{Проталкивание селекции (Predicate Pushdown):} если предикат $c$ затрагивает только атрибуты отношения $R$, то:
\begin{equation*}
\sigma_c(R \bowtie S) \;=\; \sigma_c(R) \bowtie S
\end{equation*}
Это математическое доказательство того, что фильтровать данные выгоднее \textbf{до} соединения, а не после.

\textbf{Ассоциативность соединений:}
\begin{equation*}
(R \bowtie S) \bowtie T \;=\; R \bowtie (S \bowtie T)
\end{equation*}
Это даёт оптимизатору право переставлять соединения таблиц в поиске наиболее эффективного порядка.

\vspace{0.5cm}\hrule\vspace{0.5cm}

\section{Логическая оптимизация}

Реляционная алгебра доказывает, что одинаковый результат можно получить несколькими эквивалентными способами. Задача логического оптимизатора — среди всех эквивалентных планов найти тот, который обходится дешевле всего.

Рассмотрим конкретный масштаб задачи. Предположим, таблица \texttt{employees} содержит один миллион строк, 997 000 из которых не относятся к отделу \texttt{dept\_id = 10}. Без оптимизации наивный план выглядит так:

\begin{verbatim}
Aggregate (GROUP BY name, SUM salary)
  └── Filter (dept_id = 10)
        └── Scan (employees)    ← читает все 1 000 000 строк
\end{verbatim}

Агрегатор получает миллион строк, хотя 99.7\% из них будут немедленно отфильтрованы. Оптимизатор преобразует план, применяя Predicate Pushdown:

\begin{verbatim}
Aggregate (GROUP BY name, SUM salary)
  └── Scan (employees, filter=dept_id=10)    ← читает только ~3 000 строк
\end{verbatim}

\begin{longtable}{|l|l|l|}
\hline
Метрика & Без оптимизации & После Predicate Pushdown \\ \hline
Строк до агрегатора & 1 000 000 & ~3 000 \\ \hline
Относительная экономия & — & ~99.7% \\ \hline
\end{longtable}

Второе применяемое правило — \textbf{Column Pruning}: из промежуточных результатов удаляются все колонки, не задействованные в оставшихся операциях. Для нашего запроса достаточно передавать только \texttt{name}, \texttt{salary} и \texttt{dept\_id} — все остальные атрибуты строки (если они есть) отсекаются немедленно после чтения.

В реализуемом движке оптимизатор (\texttt{planner/optimizer.py}) устроен как \textbf{набор правил (Rule-Based Optimizer, RBO)}. Каждое правило — это функция, принимающая узел логического плана и возвращающая, при выполнении условий применимости, эквивалентный узел с меньшей ожидаемой стоимостью. Правила применяются итеративно до тех пор, пока план не перестаёт меняться.

\vspace{0.5cm}\hrule\vspace{0.5cm}

\section{Модель Volcano: Итераторная модель исполнения}

После того как оптимизатор сформировал логический план, его необходимо \textbf{исполнить} — перебрать реальные данные и вычислить итоговый результат. Модель Volcano, предложенная Гётцем Грэфе в 1994 году, была стандартом реляционных СУБД на протяжении десятилетий и применяется в PostgreSQL, MySQL и множестве других систем.

Суть модели: каждый физический оператор реализует единый интерфейс итератора с методом \texttt{\textit{\_next\_}()}. Вычисление организовано по принципу \textbf{pull} — «вытягивания»: верхний оператор запрашивает следующую строку у нижележащего, тот — у своего, и так до самого нижнего оператора \texttt{Scan}, читающего данные с диска.

Проследим цепочку вызовов для нашего оптимизированного плана:

\begin{verbatim}
1.  Aggregate.__next__()
      → вызывает Scan.__next__()   ← строка 1: name="Alice", dept_id=7 → не подходит
      → вызывает Scan.__next__()   ← строка 2: name="Bob",   dept_id=10 → подходит
      → накапливает ("Bob", salary) в буфере агрегата
2.  Aggregate.__next__()
      → продолжает вызывать Scan до следующей строки с dept_id=10
      ...
N.  Scan.__next__() бросает StopIteration — данные исчерпаны
N+1 Aggregate флашит накопленные группы → возвращает финальные строки
\end{verbatim}

При таблице в миллион строк это \textbf{миллион вызовов} \texttt{Scan.\textit{\_next\_}()}. Каждый вызов — отдельный переход интерпретатора Python между кадрами стека, отдельная проверка типов, отдельное обращение к памяти.

Узкое место этой модели объясняется архитектурой современных процессоров. Процессор работает несравнимо быстрее, чем оперативная память (задержка L1-кэша — 1–4 такта, RAM — 200–300 тактов). Чтобы не простаивать, CPU должен заранее подгружать данные в кэш. При обработке одной строки (Python-словарь с атрибутами) кэш-линия заполняется лишь частично, и следующая строка снова вынуждает CPU обращаться к RAM напрямую. Возникают систематические \textbf{промахи кэша (cache misses)}, парализующие работу конвейера процессора.

Вместе с тем у итераторной модели есть неоспоримое достоинство: в любой момент времени в памяти присутствует не более одной строки данных (не считая буфера агрегата). Для транзакционных нагрузок (OLTP), где нужно быстро прочитать или записать единственную запись, Volcano остаётся оптимальным выбором.

\vspace{0.5cm}\hrule\vspace{0.5cm}

\section{Векторизованная модель исполнения}

Векторизация — это архитектурный ответ на ограничения итераторной модели. Вместо обработки одной строки за вызов оператор возвращает \textbf{пакет (batch)} строк — словарь NumPy-массивов, где каждый массив содержит данные одной колонки для тысячи строк одновременно:

\begin{verbatim}
# Volcano — один вызов, одна строка:
{"name": "Alice", "salary": 85000, "dept_id": 10}

# Vectorized — один вызов, 1024 строки:
{
  "name":    np.array(["Alice", "Bob", ..., "Zara"]),
  "salary":  np.array([85000, 92000, ..., 78000]),
  "dept_id": np.array([10, 10, ..., 10])
}
\end{verbatim}

При размере пакета 1024 строки миллион строк передаётся за \textbf{977 вызовов} вместо миллиона. Интерпретационные накладные расходы Python сокращаются пропорционально — но важно другое: \textbf{вся арифметика выполняется внутри NumPy}, написанного на C и Fortran, без динамической типизации и без GIL на уровне математических операций.

Сравним оба подхода для нашего запроса на таблице из миллиона строк:

\begin{longtable}{|l|l|l|}
\hline
Метрика & Volcano & Vectorized (batch=1024) \\ \hline
Вызовов \texttt{\textit{\_next\_}()} / \texttt{get\_next\_batch()} & 1 000 000 & ~977 \\ \hline
Фильтрация & \texttt{if row["dept\_id"] == 10} × 1M & \texttt{arr[arr["dept\_id"] == 10]} (NumPy C-код) \\ \hline
Использование L1/L2 кэша & Низкое & Высокое \\ \hline
\end{longtable}

Эффективность кэша объясняется \textbf{плотным колоночным хранением} NumPy-массивов. Когда процессор считывает первый элемент массива \texttt{salary}, он автоматически загружает в L1-кэш следующие 7 соседних значений (кэш-линия = 64 байта = 8 числа \texttt{float64}). Все последующие обращения — прямо из кэша, без ожидания RAM.

Векторизация также открывает доступ к \textbf{SIMD-инструкциям процессора} (AVX2 и аналогичным). SIMD позволяет выполнить одну операцию над несколькими парами данных за один такт: AVX2 складывает 4 пары \texttt{double} одновременно. Для нашей агрегации \texttt{SUM(salary)} это означает аппаратное параллелизование непосредственно на уровне регистров.

Принципиальное различие моделей относится к ориентации хранения данных. Volcano работает со строками (\textbf{NSM — N-ary Storage Model}): все атрибуты одной записи хранятся вместе. Векторизованный движок работает со столбцами (\textbf{DSM — Decomposition Storage Model}): при агрегации \texttt{SUM(salary)} процессор читает только колонку \texttt{salary}, не тратя пропускную способность шины памяти на \texttt{name}, \texttt{dept\_id} и любые другие атрибуты, не участвующие в текущей операции.

\vspace{0.5cm}\hrule\vspace{0.5cm}

\section{Инструментальная среда: Python и NumPy}

Выбор Python в качестве языка реализации нередко воспринимается как компромисс производительностью ради удобства разработки. Настоящий раздел обосновывает, почему для исследовательского SQL-движка этот выбор является не только приемлемым, но и методологически правильным.

Стандартный интерпретатор Python (CPython) имеет два фундаментальных ограничения. \textbf{Первое} — динамическая типизация и боксинг: каждое число Python является полноценным объектом с заголовком в 28 байт (счётчик ссылок, указатель на тип, значение). Простая операция \texttt{a + b} требует распаковки \texttt{a}, распаковки \texttt{b} и создания нового объекта-результата — это катастрофически дорого при миллионе итераций. \textbf{Второе} — Global Interpreter Lock (GIL): CPython допускает исполнение байткода только в одном потоке, что исключает настоящий параллелизм для CPU-задач.

Решение, применённое в данной работе, состоит в архитектурном разделении: \textbf{Python управляет, NumPy вычисляет}. Python-код отвечает за жизненный цикл операторов, координацию батчей и маршрутизацию данных — функции, где производительность некритична. NumPy выполняет все арифметические операции: свёртки, фильтрации, агрегации, сортировки — внутри скомпилированного C-кода, не проходя через интерпретатор и не подпадая под ограничения GIL для математических операций.

Фактически Python-интерпретатор исполняет ~977 высокоуровневых вызовов (по числу батчей), а не миллион итераций. Этот архитектурный шов — граница Python/NumPy — и есть то место, где производительность «переключается» с интерпретируемого на нативный режим.

\begin{longtable}{|l|l|l|}
\hline
Критерий & Python + NumPy & C++ / Rust \\ \hline
Скорость разработки & Высокая & Низкая \\ \hline
Скорость исполнения & ~3–15× медленнее нативного & Максимальная \\ \hline
Интроспекция и визуализация & Встроенная (Matplotlib, Graphviz) & Требует внешних инструментов \\ \hline
Читаемость исходного кода & Высокая & Средняя \\ \hline
Целевая нагрузка & OLAP-аналитика (batch) & OLTP (< 1 мс на запрос) \\ \hline
\end{longtable}

Для исследовательского движка, цель которого — продемонстрировать алгоритмические принципы и сравнить архитектурные подходы, Python является стандартным выбором в академической среде. Именно на Python построены PySpark, Apache Arrow и большинство исследовательских СУБД последних лет.

\vspace{0.5cm}\hrule\vspace{0.5cm}

\section{Визуализация и наблюдаемость системы}

Системы обработки данных, работающие как «чёрный ящик», крайне сложно отлаживать: разработчик видит входной запрос и итоговый результат, но не имеет возможности понять, что произошло внутри. В данной работе этому аспекту уделяется особое внимание — движок проектируется как \textbf{«белый ящик» (White-Box System)}.

Концепция белого ящика предполагает, что пользователь или разработчик в любой момент может наблюдать: как оптимизатор интерпретировал запрос (физический план), сколько строк прошло через каждый оператор, и где именно тратится большая часть времени.

Для нашего демонстрационного запроса движок генерирует следующий физический план в виде ориентированного ациклического графа (DAG):

\begin{verbatim}
[CSV: employees.csv]
         │  1 000 000 строк
         ▼
[FilterOp: dept_id = 10]
         │  ~3 000 строк        ←  здесь виден эффект Predicate Pushdown
         ▼
[AggregateOp: GROUP BY name, SUM(salary)]
         │  N уникальных групп
         ▼
[SortOp: ORDER BY SUM(salary) DESC]
         │  N строк
         ▼
[ProjectOp: SELECT name, SUM(salary)]
         │  N строк
         ▼
      РЕЗУЛЬТАТ
\end{verbatim}

Числа на рёбрах делают эффект оптимизации \textbf{немедленно очевидным визуально}: переход от миллиона строк до трёх тысяч виден на диаграмме даже без чтения кода. Такой граф сохраняется движком в директорию \texttt{bench\_output/} и доступен как Graphviz-изображение или ASCII-дерево.

Дополнительным инструментом является \textbf{тепловая карта производительности}: узлы DAG окрашиваются пропорционально их доле в суммарном времени выполнения — от синего (незначительные затраты) до красного (узкое место). Это превращает оптимизацию из слепого перебора вариантов в управляемый процесс: разработчик видит, какой именно оператор необходимо улучшить, не анализируя профайлинговые логи вручную.

\vspace{0.5cm}\hrule\vspace{0.5cm}

\section{Выводы по первой главе}

Первая глава была посвящена теоретическому фундаменту, без которого проектирование SQL-движка превращается в набор эмпирических решений без понимания их причин. Ключевая методическая идея — прослеживание единого запроса через все этапы обработки — позволила показать, что каждый теоретический концепт непосредственно соответствует конкретному инженерному решению.

Маршрут запроса выглядит следующим образом:

\begin{verbatim}
«SELECT name, SUM(salary) FROM employees WHERE dept_id=10 GROUP BY name ORDER BY SUM(salary) DESC»
       │
       ├─[Лексер]──────────────────────► 24 токена
       │
       ├─[Парсер LALR(1)]──────────────► AST: SelectNode(...)
       │
       ├─[Логический планировщик]──────► σ ∘ γ ∘ τ ∘ π (реляционная алгебра)
       │
       ├─[Оптимизатор RBO]─────────────► Predicate Pushdown: 1M строк → 3K
       │
       ├─[Volcano / Vectorized]────────► физическое исполнение
       │
       └─[Визуализация DAG]────────────► наблюдаемость результата
\end{verbatim}

Связь между теоретическими концепциями и программными модулями реализации, описываемыми в Главе 2, представлена в следующей таблице:

\begin{longtable}{|l|l|l|}
\hline
Теоретическая концепция (Гл. 1) & Модуль реализации (Гл. 2) & Файл проекта \\ \hline
Декларативность / DSL (§1.1) & Архитектура движка (§2.1) & \texttt{README.md} \\ \hline
Формальная грамматика CFG (§1.2) & Грамматика Lark (§2.2) & \texttt{parser/grammar.lark} \\ \hline
Лексический анализ (§1.3) & Лексер и терминалы (§2.2) & \texttt{parser/grammar.lark} \\ \hline
LALR-парсер / AST (§1.4) & Построитель AST (§2.3) & \texttt{parser/ast\_builder.py} \\ \hline
Реляционная алгебра (§1.5) & Логические операторы (§2.6) & \texttt{planner/logical\_operators.py} \\ \hline
Predicate Pushdown / RBO (§1.6) & Правила оптимизатора (§2.7) & \texttt{planner/optimizer.py} \\ \hline
Модель Volcano (§1.7) & Итераторный исполнитель (§2.8) & \texttt{executor/volcano\_executor.py} \\ \hline
Векторизованная модель (§1.8) & Векторный исполнитель (§2.9) & \texttt{executor/vectorized\_executor.py} \\ \hline
Python + NumPy (§1.9) & Инструментарий проекта (§2.1) & \texttt{requirements.txt} \\ \hline
White-box / визуализация (§1.10) & Инструменты визуализации (§2.11) & \texttt{visualization/}, \texttt{bench\_output/} \\ \hline
\end{longtable}

Из проведённого анализа можно сформулировать несколько принципиальных выводов.

Декларативность SQL — не синтаксический сахар, а фундаментальная архитектурная идея, разделяющая описание задачи и способ её решения. Именно это разделение делает возможной автоматическую оптимизацию.

Принадлежность SQL к классу контекстно-свободных языков (Тип 2 по Хомскому) — это не историческая случайность, а следствие требования поддерживать произвольно вложенные выражения. Это требование однозначно диктует использование LALR-парсера и стекового автомата.

Законы реляционной алгебры — это не эвристики и не «хорошие практики»: это математические теоремы. Правило Predicate Pushdown представляет собой доказуемое утверждение об эквивалентности двух выражений, что гарантирует корректность оптимизатора независимо от конкретных данных.

Разница между моделями Volcano и Vectorized проявляется не в правильности результата, а в эффективности использования аппаратных ресурсов. Выбор между ними — инженерный компромисс между простотой реализации (Volcano) и производительностью на аналитической нагрузке (Vectorized).

Изложенная теория является непосредственным основанием для всех архитектурных решений, детально описанных в Главе 2.

\end{document}
